
\documentclass[12pt,a4paper]{article}

\usepackage[a4paper,margin=1in]{geometry}
\usepackage{graphicx}
\usepackage{booktabs}
\usepackage{amsmath,amssymb}
\usepackage{float}
\usepackage{hyperref}
\usepackage{xcolor}
\usepackage{listings}
\usepackage{enumitem}
\usepackage{fancyhdr}
\usepackage{longtable}

\hypersetup{colorlinks=true,linkcolor=blue,urlcolor=blue}

\pagestyle{fancy}
\fancyhf{}
\lhead{ICS1512}
\rhead{Machine Learning Laboratory}
\cfoot{\thepage}

\lstset{
language=Python,
basicstyle=\ttfamily\small,
numbers=left,
frame=single,
breaklines=true,
showstringspaces=false
}

\begin{document}

\begin{tabular}{|p{4cm}|p{11cm}|}
\hline
Experiment No. & 7 \\ \hline
Experiment Title & Bagging, Boosting, and Stacked Ensemble Models
\footnote{\textbf{GitHub Repository: }\url{https://github.com/Barath-j593/ML-assignments}}\\ \hline

Student Name & NavinKumar S \\ \hline
Register Number & 3122247001039 \\ \hline
Faculty & Dr. Poreddy Ajay Kumar Reddy \\ \hline
Submission Date & 02/08/26 \\ \hline
\end{tabular}

\section{Objective}
To understand ensemble learning strategies -- Bagging, Boosting, and Stacking -- implement
Bagging and Boosting classifiers, build a Stacked Ensemble using multiple heterogeneous base
learners, compare all three approaches in terms of accuracy, stability, and generalization,
and analyze their effect on the bias--variance trade-off.

\section{Problem Statement}
Given the Wisconsin Diagnostic Breast Cancer dataset of 30 numerical features, the problem is
binary classification: predict malignant vs. benign. The objective is to determine whether
combining multiple models -- via bootstrap aggregation (Bagging), sequential error correction
(Boosting), or a learned meta-combination (Stacking) -- improves accuracy and stability over
individual base learners, and to characterize how each ensemble strategy affects bias and
variance differently.

\section{Dataset Description}
\begin{longtable}{|p{6cm}|p{8cm}|}
\hline
Dataset Name & Wisconsin Diagnostic Breast Cancer (WDBC) \\ \hline
Dataset Source & UCI Machine Learning Repository -- \url{https://archive.ics.uci.edu} (loaded via \texttt{sklearn.datasets.load\_breast\_cancer}) \\ \hline
Number of Samples & 569 \\ \hline
Number of Features & 30 numerical attributes (mean, standard error, and ``worst'' value of 10 cell-nucleus measurements) \\ \hline
Number of Classes & 2 (357 benign / 212 malignant) \\ \hline
Missing Values & 0 \\ \hline
Train-Test Split & 80:20 stratified (455 train / 114 test), \texttt{random\_state=42} \\ \hline
\end{longtable}

\section{Exploratory Data Analysis}
Include all relevant visualizations.
\begin{itemize}
\item Dataset overview
\item Statistical summary
\item Missing value analysis
\item Class distribution
\item Correlation matrix
\item Feature distribution
\end{itemize}

\subsection*{Class Distribution}
\begin{figure}[H]
\centering
\includegraphics[width=0.5\linewidth]{class_distribution.png}
\caption{Class distribution -- Breast Cancer Wisconsin (Diagnostic)}
\end{figure}
\textbf{Inference:} The dataset is mildly imbalanced (357 benign vs. 212 malignant, about
63:37). No missing values are present, so no missing-value visualization is required.

\subsection*{Correlation Matrix}
\begin{figure}[H]
\centering
\includegraphics[width=0.85\linewidth]{correlation_heatmap.png}
\caption{Correlation heatmap -- top 12 features correlated with diagnosis}
\end{figure}
\textbf{Inference:} ``Worst concave points'', ``worst perimeter'', ``worst radius'', and
``mean concave points'' correlate most strongly with diagnosis ($|r| > 0.7$), the same
features that dominated feature importance in the single Decision Tree / Random Forest
experiment on this dataset -- these cell-nucleus size and shape irregularities are the
primary drivers every model in this experiment learns from.

\section{Data Preprocessing}
Explain:
\begin{itemize}
\item \textbf{Label encoding:} The target is already numerically encoded by scikit-learn
(0 = malignant, 1 = benign); no additional encoding step was required.
\item \textbf{Missing value handling:} Not required -- the dataset contains zero missing
values across all 569 rows and 30 features.
\item \textbf{Feature scaling:} Not required for the Decision Tree, Bagging, AdaBoost, or
Gradient Boosting components (tree-based, invariant to monotonic scaling). The Stacked
Ensemble's SVM base learner is technically scale-sensitive, but was left unscaled here
consistent with the assignment's base-model specification; this is noted as a limitation in
Section 14.
\item \textbf{Train-test split:} 80:20 stratified split (455 train / 114 test rows),
preserving the 63:37 class ratio in both subsets.
\end{itemize}

\section{Machine Learning Algorithm}
Explain:
\begin{itemize}
\item \textbf{Working principle:} \textbf{Bagging} trains many independent copies of a base
estimator (here, Decision Trees) on different bootstrap samples of the training data and
aggregates their predictions by majority vote -- primarily reduces variance. \textbf{Boosting}
trains models sequentially, where each new model focuses on the samples the previous ones
got wrong (AdaBoost re-weights misclassified samples; Gradient Boosting fits new trees to the
residual errors of the ensemble so far) -- primarily reduces bias. \textbf{Stacking} trains
several diverse base models (here: SVM, Gaussian Naive Bayes, Decision Tree) and a
meta-learner (Logistic Regression) that learns how to optimally combine their predictions,
using cross-validated out-of-fold predictions as the meta-learner's training input.
\item \textbf{Mathematical formulation:}
$$\hat{y}_{\text{bag}} = \operatorname*{mode}_{i=1,\dots,B} T_i(x), \quad T_i \text{ trained on bootstrap sample } i$$
$$F_m(x) = F_{m-1}(x) + \eta \cdot h_m(x) \quad \text{(Gradient Boosting, } \eta \text{ = learning rate)}$$
$$\hat{y}_{\text{stack}} = g\big(f_1(x), f_2(x), \dots, f_k(x)\big) \quad \text{(} g \text{ = meta-learner, } f_i \text{ = base learners)}$$
\item \textbf{Advantages:} Bagging is simple, parallelizable, and reliably reduces variance
for high-variance base learners like deep trees. Boosting can achieve very high accuracy by
directly targeting the errors of previous models. Stacking can outperform any single
ensemble strategy by exploiting complementary strengths of heterogeneous models.
\item \textbf{Limitations:} Bagging cannot reduce bias -- if the base learner is
systematically wrong, averaging many copies of it will not fix that. Boosting is sequential
(cannot parallelize across estimators) and more prone to overfitting on noisy data if
learning rate/depth are not controlled. Stacking is the most computationally expensive (it
internally cross-validates every base learner) and, as observed in this experiment
(Section 12), can be less stable across folds than simpler ensembles if the meta-learner
overfits to base-learner quirks in a small dataset.
\end{itemize}

\section{Experimental Setup}
\begin{tabular}{ll}
\toprule
Python Version & \\
Scikit-Learn Version & \\
IDE & Jupyter Notebook \\
Operating System & \\
Hardware & \\
\bottomrule
\end{tabular}

\section{Hyperparameter Selection}

\begin{longtable}{|p{2.6cm}|p{2.6cm}|p{5.7cm}|p{2.3cm}|}
\hline
\textbf{Model} & \textbf{Search Method} & \textbf{Best Parameters} & \textbf{Best CV Accuracy} \\ \hline
\endhead
Bagging & Grid Search & n\_estimators=50, max\_samples=0.5, max\_features=0.5 & 0.9604 \\ \hline
Boosting (Gradient Boosting) & Grid Search & n\_estimators=200, learning\_rate=0.1, max\_depth=1 & 0.9758 \\ \hline
AdaBoost (secondary) & Grid Search & n\_estimators=100, learning\_rate=1.0 & 0.9802 \\ \hline
Stacked Ensemble & Base model choice (SVM+NB+DT vs. SVM+DT) & SVM + Gaussian NB + Decision Tree $\rightarrow$ Logistic Regression & 0.9538 \\ \hline
\end{longtable}

\textbf{Note:} Bagging was searched over n\_estimators $\in \{10,50,100,200\} \times
max\_samples \in \{0.5,0.7,1.0\} \times max\_features \in \{0.5,0.7,1.0\}$ (36 combinations).
Boosting (Gradient Boosting, used as the primary tunable representative) was searched over
n\_estimators $\in \{50,100,200\} \times$ learning\_rate $\in \{0.01,0.1,0.5,1.0\} \times
max\_depth \in \{1,2,3\}$ (36 combinations); AdaBoost was separately tuned over n\_estimators
and learning\_rate as a secondary boosting comparison. For Stacking, the assignment's
hyperparameter space is the choice of base models and final estimator rather than numeric
values: the specified (SVM, Naive Bayes, Decision Tree $\to$ Logistic Regression) combination
achieved 0.9538 average CV accuracy, versus 0.9231 for a reduced (SVM, Decision Tree)
combination -- confirming that the third, heterogeneous base learner (Naive Bayes)
meaningfully helped the meta-learner.

\subsection*{Bagging Cross-Fold Results (summarized by n\_estimators, max\_samples)}
\begin{longtable}{|p{3cm}|p{3cm}|p{4cm}|}
\hline
\textbf{n\_estimators} & \textbf{max\_samples} & \textbf{Avg CV Accuracy} \\ \hline
\endhead
10 & 1.0 & 0.9538 \\ \hline
50 & 0.7 & 0.9553 \\ \hline
50 & 1.0 & 0.9553 \\ \hline
100 & 0.7 & 0.9590 \\ \hline
100 & 1.0 & 0.9560 \\ \hline
200 & 0.7 & 0.9590 \\ \hline
200 & 1.0 & 0.9568 \\ \hline
\end{longtable}
\textbf{Inference:} Accuracy is fairly stable (0.950--0.959) across the grid, with a mild
peak around n\_estimators=100--200 and max\_samples=0.7 -- Bagging shows much less
hyperparameter sensitivity than the Decision Tree's own depth parameter did in the earlier
single-tree experiment.

\subsection*{Boosting Cross-Fold Results (summarized by n\_estimators, learning\_rate)}
\begin{longtable}{|p{3cm}|p{3cm}|p{4cm}|}
\hline
\textbf{n\_estimators} & \textbf{learning\_rate} & \textbf{Avg CV Accuracy} \\ \hline
\endhead
50 & 0.01 & 0.9238 \\ \hline
50 & 0.5 & 0.9641 \\ \hline
100 & 0.1 & 0.9575 \\ \hline
100 & 0.5 & \textbf{0.9685} \\ \hline
200 & 0.1 & 0.9678 \\ \hline
200 & 0.5 & 0.9685 \\ \hline
200 & 1.0 & 0.9634 \\ \hline
\end{longtable}
\textbf{Inference:} Very low learning\_rate (0.01) clearly underperforms regardless of
n\_estimators (0.924 at 50 estimators) -- too little correction is applied per round. Accuracy
is best and fairly flat across moderate-to-high learning rates (0.5--1.0) once enough
estimators are used, showing learning\_rate and n\_estimators trade off against each other
(fewer, larger steps vs. more, smaller steps reaching a similar destination).

\section{Results}

\begin{longtable}{|p{4cm}|p{2cm}|p{1.8cm}|p{1.8cm}|p{1.8cm}|}
\hline
\textbf{Model} & \textbf{Accuracy} & \textbf{Precision} & \textbf{Recall} & \textbf{F1 Score} \\ \hline
\endhead
Bagging & 0.9561 & 0.9589 & 0.9722 & 0.9655 \\ \hline
Boosting (Gradient Boosting) & 0.9561 & 0.9467 & 0.9861 & 0.9660 \\ \hline
Stacked Ensemble & 0.9561 & 0.9467 & 0.9861 & 0.9660 \\ \hline
\end{longtable}
\textbf{Note:} All three tuned ensembles reach identical test-set Accuracy (0.9561) on this
114-sample held-out test set, with Boosting and Stacking producing the exact same predictions
on this split (identical Precision/Recall/F1) -- a reflection of how well-separated this
dataset is, where several strong models converge on the same decisions for most test points.
Section 12's cross-validation results (which use many different splits) are the more
reliable basis for distinguishing between them.

\section{Result Visualization}

Include all graphs generated during the experiment.

\begin{figure}[H]
\centering
\includegraphics[width=0.95\linewidth]{bias_variance_curves.png}
\caption{Bias--variance illustration: train/test accuracy vs. n\_estimators for Bagging (left) and Boosting (right)}
\end{figure}
\textbf{Inference:} Both ensembles reach $\approx$100\% training accuracy almost immediately
and stay there, but their \emph{test} accuracy behaves very differently in character: Bagging's
test accuracy rises quickly and then stays essentially flat/stable regardless of how many
trees are added (0.921 $\to$ 0.948, no decline) -- consistent with variance reduction, where
more independent estimators simply average out noise without adding new capacity. Boosting's
test accuracy dips slightly around 10--25 estimators before climbing to its own stable plateau
($\approx$0.956) -- consistent with it actively reducing bias through sequential correction,
which can temporarily overshoot on limited data before stabilizing.

\begin{figure}[H]
\centering
\includegraphics[width=0.75\linewidth]{confusion_matrices.png}
\caption{Confusion matrices -- Bagging, Boosting, Stacking}
\end{figure}
\textbf{Inference:} Boosting and Stacking produce identical confusion matrices on this test
split (both misclassify the same single malignant case as benign), while Bagging makes one
additional error of the opposite type -- consistent with their matching Precision/Recall/F1
figures in Section 9.

\begin{figure}[H]
\centering
\includegraphics[width=0.55\linewidth]{roc_curves.png}
\caption{ROC curves -- Bagging vs. Boosting vs. Stacking}
\end{figure}
\textbf{Inference:} Boosting achieves the highest AUC (0.995), narrowly ahead of Bagging
(0.993) and Stacking (0.986) -- all three rank predictions very well, with Boosting's
sequential error-correction giving it a slight edge in separating the classes across
thresholds, not just at the default 0.5 cutoff.

\begin{figure}[H]
\centering
\includegraphics[width=0.7\linewidth]{cv_accuracy.png}
\caption{5-fold cross-validation accuracy -- Bagging vs. Boosting vs. Stacking}
\end{figure}
\textbf{Inference:} This is the most informative comparison in the experiment: Boosting has
the best average (0.9578) and Bagging is close behind (0.9508), but \textbf{Stacking has both
the lowest average (0.9403) and by far the largest fold-to-fold variance}, dropping to 0.895
on fold 2 while Bagging and Boosting stay above 0.93 on every fold. This directly contradicts
the naive expectation that "more sophisticated ensembling always wins" -- on this
569-sample dataset, Stacking's extra meta-learning layer adds instability rather than benefit.

\begin{figure}[H]
\centering
\includegraphics[width=0.7\linewidth]{classifier_comparison.png}
\caption{Ensemble model comparison -- Accuracy / Precision / Recall / F1}
\end{figure}
\textbf{Inference:} On the single test split all three appear virtually tied, which is why
the cross-validation view (Figure~5) -- not this single-split comparison -- should drive the
final model choice.

\section{Observation Questions}
\begin{enumerate}
\item \textbf{How does Bagging reduce variance?} By training many Decision Trees
independently on different bootstrap samples and majority-voting their predictions, the
random, sample-specific errors of individual trees tend to cancel out in the aggregate, while
genuine signal (present in most bootstrap samples) is reinforced. This is directly visible in
Figure~4 (left): test accuracy stabilizes rather than degrading as more trees are added,
unlike a single deep tree's overfitting behaviour.
\item \textbf{How does Boosting address model bias?} Each new estimator in the sequence is
trained specifically to correct the residual errors of the ensemble built so far (Gradient
Boosting fits to the gradient of the loss; AdaBoost up-weights previously misclassified
samples). This directly targets systematic (biased) errors rather than just averaging away
noise, which is why Boosting achieved the highest tuned CV accuracy of any single ensemble
(AdaBoost: 0.9802, Gradient Boosting: 0.9758) in this experiment.
\item \textbf{Why does stacking benefit from heterogeneous models?} Different model families
make different kinds of errors: SVM captures margin-based non-linear structure, Naive Bayes
captures probabilistic feature-independence patterns, and Decision Tree captures axis-aligned
rule-based structure. A meta-learner can exploit disagreements between these different error
profiles to make better combined decisions than any one model alone -- confirmed here by the
alternative-combination check (Section 8): removing Naive Bayes from the base learners dropped
CV accuracy from 0.9538 to 0.9231, showing the heterogeneity genuinely mattered.
\item \textbf{Which ensemble method performed best and why?} By 5-fold cross-validation (the
most reliable comparison, Section 12), \textbf{Boosting} performed best (avg. 0.9578),
narrowly ahead of Bagging (0.9508), with \textbf{Stacking performing worst and least stable}
(avg. 0.9403, largest fold variance). Boosting's sequential bias-correction was well suited to
this already fairly separable dataset; Stacking's added complexity (an extra meta-learning
layer trained on limited data) appears to have introduced instability rather than benefit at
this dataset size.
\end{enumerate}

\section{Discussion}
This experiment demonstrates that ensemble methods are not interchangeable, "always better"
upgrades over single models -- each targets a different failure mode. Bagging's stable,
flat test-accuracy curve as n\_estimators grows (Figure~4) is a clean illustration of pure
variance reduction. Boosting achieved the best overall cross-validated accuracy, consistent
with its bias-reduction mechanism directly attacking the (already fairly small) residual
error left after any single tree. Stacking, despite combining three heterogeneous and
individually reasonable base learners, was the \emph{least} stable ensemble by cross-validation
-- a useful, non-obvious finding: on a relatively small dataset (569 samples, 455 training
rows split further across the meta-learner's internal cross-validation), a stacked
meta-learner has limited data to learn a reliable combination rule, and can end up more
sensitive to which samples land in which fold than simpler ensembles. A limitation of this
experiment is that the Stacked Ensemble's SVM base learner was not feature-scaled, which
likely handicapped its individual contribution and, by extension, the stack's overall
stability -- a natural follow-up would be to scale features before stacking and observe
whether that narrows the gap to Boosting.

\section{Conclusion}
For the Wisconsin Diagnostic Breast Cancer dataset, Boosting (Gradient Boosting, tuned to
n\_estimators=200, learning\_rate=0.1, max\_depth=1, or equivalently AdaBoost at
n\_estimators=100, learning\_rate=1.0) is the best-performing and most reliable ensemble
strategy by 5-fold cross-validation, narrowly ahead of Bagging. Stacking, despite its
theoretical appeal of combining heterogeneous learners, was both the weakest and least stable
ensemble on this dataset -- a reminder that ensemble sophistication does not guarantee better
generalization, particularly on smaller datasets where a meta-learner has limited data to
learn a robust combination rule. All three ensembles nonetheless comfortably improved on a
single unconstrained Decision Tree's known overfitting behaviour (established in the prior
single-tree experiment on this same dataset), confirming the core value of ensembling even
where the specific strategy matters.

\section{References}
\begin{enumerate}[label={[\arabic*]}]
\item Scikit-learn Documentation, ``Ensemble Methods,'' \url{https://scikit-learn.org/stable/modules/ensemble.html}
\item Scikit-learn Documentation, ``Bagging Classifier,'' \url{https://scikit-learn.org/stable/modules/generated/sklearn.ensemble.BaggingClassifier.html}
\item UCI Machine Learning Repository, ``Breast Cancer Wisconsin (Diagnostic) Data Set,'' \url{https://archive.ics.uci.edu}
\end{enumerate}

\end{document}
